\chapter{Conception de l'approche HAKE-MER}
\label{chap:conception}

\section*{Introduction}
\addcontentsline{toc}{section}{Introduction}
Le chapitre~\ref{chap:etat_art} a fixé le corpus GoEmotions, une baseline
par modèles pré-entraînés et les réponses littéraires aux limitations
L1 à L3. Ce chapitre en déduit l'architecture \textbf{HAKE-MER}
(\emph{Hierarchical Attention with Knowledge Enhancement for
Multi-Emotion Recognition}) : formalisation de la tâche, description
des modules M1 à M4 et hypothèses H1 à H4 que le
chapitre~\ref{chap:experimentation} fera trancher par ablation.

\section{Formulation du problème et limites existantes}
\label{sec:conception-formulation}

Soit un texte $x$ et un inventaire d'émotions
$\mathcal{E}=\{e_1,\dots,e_K\}$ ($K{=}28$ pour GoEmotions).
La classification multi-étiquettes apprend une
fonction $f : x \mapsto Y \subseteq \mathcal{E}$, réalisée par $K$
sorties sigmoïdes $\hat{y}_k = \sigma(z_k)$ entraînées par entropie
croisée binaire, variante de la \emph{binary relevance}
\cite{zhang2014multilabel} (une perte par étiquette, sans modèle
explicite des co-occurrences) :
\begin{equation}
\mathcal{L}_{\mathrm{BCE}}
= -\frac{1}{K}\sum_{k=1}^{K}
\left[\, y_k \log \hat{y}_k + (1-y_k)\log(1-\hat{y}_k) \,\right].
\label{eq:bce}
\end{equation}
Une étiquette est prédite si $\hat{y}_k > 0{,}5$.

Trois limites des approches standards motivent la conception :

\begin{itemize}
\item \textbf{L1 : Indépendance des étiquettes.} La tête
  « pertinence binaire » sur un vecteur global partagé ignore les
  co-occurrences d'émotions, pourtant systématiques dans les données
  \cite{demszky2020} ;
\item \textbf{L2 : Cécité à la structure du texte.} Une
  représentation plate (le seul vecteur \texttt{[CLS]} d'un encodeur
  type BERT \cite{devlin2019}) mélange les signaux portés par des
  phrases différentes ; les commentaires Reddit
  de GoEmotions comportent souvent plusieurs phrases \cite{demszky2020} ;
\item \textbf{L3 : Pauvreté du signal pour les émotions rares.} Les
  fréquences d'émotions restent très inégales (par exemple l'admiration
  environ trente fois plus fréquente que le deuil sur le corpus)
  \cite{demszky2020}, ce qui limite l'apprentissage purement supervisé
  sur la longue traîne.
\end{itemize}

\section{Vue d'ensemble de l'architecture}
\label{sec:conception-vue}

HAKE-MER répond à ces limites par quatre modules composables,
construits au-dessus d'un backbone pré-entraîné partagé
(RoBERTa, DistilBERT ou BERT \cite{devlin2019,liu2019roberta,sanh2019distilbert},
figure~\ref{fig:hakemer-archi}) : le \textbf{Module 1 (Encodage
Hiérarchique)} préserve la structure du texte (L2) ; le
\textbf{Module 2 (Attention Croisée Multi-Émotions)} modélise les
interactions entre émotions (L1) ; le \textbf{Module 3
(Enrichissement Sémantique)} injecte une connaissance affective
externe (L3) ; le \textbf{Module 4 (Pondération Dynamique)} adapte
l'intensité de chaque émotion au contexte. Chaque module s'active par
un indicateur de configuration, ce qui permet l'ablation cumulative
du chapitre 4.

\begin{figure}[H]
\centering
\resizebox{\textwidth}{!}{%
\begin{tikzpicture}[
  font=\small,
  node distance=4mm and 6mm,
  box/.style={draw, rounded corners=2pt, align=center, minimum height=9mm, inner sep=3pt},
  mod/.style={box, thick},
  flow/.style={-{Stealth[length=2.5mm]}, thick}
]
\node[box, fill=gray!12] (input) {Texte d'entrée\\ $M$ phrases $\times$ 32 tokens};
\node[box, fill=gray!22, right=of input] (plm) {Backbone partagé\\ (RoBERTa / DistilBERT)};
\node[mod, fill=red!14, right=of plm] (m1) {\textbf{Module 1}\\ Encodage Hiérarchique\\ attention mot$\to$phrase};
\node[mod, fill=blue!12, right=of m1] (xatt) {\textbf{Module 2}\\ attention croisée\\ émotion$\to$phrases};
\node[box, fill=blue!8, right=of xatt] (inter) {auto-attention\\ inter-émotions};
\node[mod, fill=violet!14, right=of inter] (m4) {\textbf{Module 4}\\ Pondération\\ Dynamique};
\node[box, fill=gray!25, right=of m4] (out) {$K$ sorties\\ sigmoïdes};
\node[box, fill=blue!6, below=9mm of xatt] (q) {$K$ requêtes d'émotions\\ $Q_k + g\cdot\mathrm{prior}$};
\node[mod, fill=green!16, below=9mm of q] (gate) {\textbf{Module 3}\\ Enrichissement Sémantique\\ porte $g$ (init.~0)};
\node[box, fill=green!8, left=of gate] (proj) {projection\\ $+\;\tanh$};
\node[box, fill=green!12, left=of proj] (nrc) {Lexique affectif\\ (NRC / SenticNet)};
\draw[flow] (input) -- (plm);
\draw[flow] (plm) -- (m1);
\draw[flow] (m1) -- node[above, font=\footnotesize\itshape] {$s_1..s_M$} (xatt);
\draw[flow] (xatt) -- (inter);
\draw[flow] (inter) -- (m4);
\draw[flow] (m4) -- (out);
\draw[flow] (q) -- (xatt);
\draw[flow] (nrc) -- (proj);
\draw[flow] (proj) -- (gate);
\draw[flow] (gate) -- (q);
\end{tikzpicture}}
\caption{Architecture de HAKE-MER : quatre modules composables
au-dessus d'un backbone partagé.}
\label{fig:hakemer-archi}
\end{figure}

\section{Module 1 : Encodage Hiérarchique}
\label{sec:conception-m1}

Le texte est segmenté en phrases ($M{=}4$ au plus, 32 tokens par
phrase, bornes fixées au chapitre~\ref{chap:experimentation}), encodées par
le backbone en un seul appel vectorisé. Le module s'inspire de
l'agrégation mot$\to$phrase des réseaux d'attention hiérarchiques
\cite{yang2016}, avec une couche d'\emph{attention additive}
\cite{bahdanau2015}. Pour une phrase de représentations de tokens
$H=(h_1,\dots,h_T)$ et de masque $m$, cette couche calcule :
\begin{equation}
u_t = v^\top \tanh(W h_t), \qquad
\alpha_t = \frac{\exp(u_t)}{\sum_{t'} \exp(u_{t'})}, \qquad
s = \sum_t \alpha_t\, h_t,
\label{eq:attnpool}
\end{equation}
les positions masquées recevant un score $-\infty$ avant le softmax.
On obtient les vecteurs de phrases $s_1,\dots,s_M$. Utilisé seul, le
module applique une seconde attention additive phrase$\to$document et
une tête linéaire partagée ; au sein du modèle complet, ce sont les
vecteurs de phrases qui alimentent le module 2.

\section{Module 2 : Attention Croisée Multi-Émotions}
\label{sec:conception-m2}

Chaque émotion $k$ possède une requête apprise $Q_k \in \mathbb{R}^d$
(initialisée en $\mathcal{N}(0, 0{,}02^2)$). L'attention croisée
multi-têtes (8 têtes, mécanisme Transformer \cite{vaswani2017}) utilise
les requêtes d'émotions comme \emph{queries} et les vecteurs de
phrases comme \emph{keys/values} :
\begin{equation}
r_k = \mathrm{MHA}\big(Q_k,\; [s_1,\dots,s_M],\; [s_1,\dots,s_M]\big),
\label{eq:mha}
\end{equation}
produisant une représentation contextuelle \emph{par émotion} :
chaque émotion « lit » le texte à sa manière. Une couche d'encodeur Transformer \cite{vaswani2017} appliquée sur
l'axe des émotions modélise ensuite leurs interactions (renforcement
ou inhibition mutuels), en complément des approches par chaînes de
classificateurs \cite{read2009classifier}. La tête de classification
est \emph{par étiquette} : chaque émotion a son propre vecteur de
poids, $z_k = w_k^\top r_k + b_k$.

\section{Module 3 : Enrichissement Sémantique}
\label{sec:conception-m3}

Un vecteur affectif de document $c \in \mathbb{R}^{d_c}$ est extrait
d'une ressource externe. Pour le \textbf{NRC Emotion Lexicon}, on
agrège au niveau document un score par émotion de base ($d_c{=}8$,
cadre Plutchik \cite{mohammad2013}), puis on applique une table de
correspondance vers les $K{=}28$ étiquettes GoEmotions (détail du
mapping au chapitre~\ref{chap:experimentation}). Pour
\textbf{SenticNet}, on agrège des scores conceptuels continus issus de
SenticNet~5 \cite{cambria2018senticnet} (primitives et polarité
conceptuelle), regroupés en un vecteur $c$ de dimension $d_c$ fixée au
protocole (chapitre~\ref{chap:experimentation}). Dans les deux cas,
$c$ est projeté, borné, puis injecté dans
les requêtes d'émotions via une \emph{porte apprise} $g$ initialisée à
zéro :
\begin{equation}
\mathrm{prior} = \tanh(W_c\, c), \qquad
\tilde{Q}_k = Q_k + g \cdot \mathrm{prior}.
\label{eq:gate}
\end{equation}
À l'initialisation ($g{=}0$), le module 3 est strictement équivalent
au module 2 : le modèle apprend la quantité de connaissance à
injecter, et la valeur finale $|g|$ mesure directement l'utilisation
de la ressource, une sonde d'interprétabilité intégrée.

\section{Module 4 : Pondération Dynamique}
\label{sec:conception-m4}

Une porte d'intensité par émotion, conditionnée par le contexte
émotionnel global $\bar{r} = \frac{1}{K}\sum_k r_k$ :
\begin{equation}
\gamma_k = \sigma\!\big(W_2\,\mathrm{ReLU}(W_1 [r_k ; \bar{r}]) + b_2\big),
\qquad r_k \leftarrow \gamma_k \cdot r_k,
\label{eq:dyngate}
\end{equation}
avec $b_2$ initialisé à $+3$ (soit $\gamma_k \approx 0{,}95$) : le
module démarre proche de l'identité et apprend à moduler
progressivement l'intensité de chaque émotion selon le texte.

\section{Hypothèses de recherche}
\label{sec:conception-hypotheses}

Le protocole du chapitre~\ref{chap:experimentation} est conçu pour
tester quatre hypothèses. Chacune lie un module à une prédiction
vérifiable sur GoEmotions, sans supposer de gains déjà mesurés :

\begin{itemize}
\item[\textbf{H1.}] La hiérarchie explicite (M1) améliore la
  reconnaissance des émotions rares (F1-macro) en préservant les
  signaux par phrase ;
\item[\textbf{H2.}] La modélisation des interactions entre émotions
  (M2) apporte un gain robuste par rapport à la prédiction
  indépendante ;
\item[\textbf{H3.}] L'injection d'un lexique affectif (M3) aide là où
  la supervision manque, et une ressource \emph{catégorielle} alignée
  sur la taxonomie cible (NRC) est plus efficace qu'une ressource
  \emph{dimensionnelle} abstraite (SenticNet) ;
\item[\textbf{H4.}] Chaque module doit justifier son ajout par
  ablation contrôlée : la plausibilité architecturale ne suffit pas.
\end{itemize}

\section*{Conclusion}
\addcontentsline{toc}{section}{Conclusion}
Ce chapitre a présenté la conception de HAKE-MER : quatre modules
composables répondant aux limitations L1 à L3, avec deux portes apprises
et interprétables (connaissance externe, intensité). Le
chapitre~\ref{chap:experimentation} précise le protocole d'ablation à
budget constant, les backbones et les analyses qui permettront de
valider ou d'infirmer H1 à H4.
